\documentclass[10pt,a4paper]{article}
\usepackage{amsmath,amssymb,amsthm,graphicx}
\usepackage{amsmath,amssymb,amsthm,hyperref}
\newtheorem{nttheorem}{Theorem}
\newtheorem{ntdefinition}[nttheorem]{Definition}
\newtheorem{ntproposition}[nttheorem]{Proposition}
\newtheorem{ntlemma}[nttheorem]{Lemma}
\newtheorem{ntcorollary}[nttheorem]{Corollary}
\title{Limits of separate preprocessing for inner product and constant-depth simultaneous-message referees}
\author{Yuval Filmus, Yuval Ishai, Avi Kaplan, Guy Kindler}
\date{doi:10.4230/LIPIcs.CCC.2020.17; CC BY3.0}
\begin{document}
\maketitle
\noindent\textit{Editorial note.} Yuval Filmus, Yuval Ishai, Avi Kaplan and Guy Kindler. Limits of Preprocessing. CCC2020 article17,DOI10.4230/LIPIcs.CCC.2020.17. CC BY3.0.\par\medskip
\noindent\textit{Editorial note.} Faithful editable transcription of the official CCC2020 CC BY3.0 publication, article17 pp17:1--17:22. One main preprocessing lower-bound mechanism: original formal Theorem9, full Section2 definitions/standard prerequisites3--8, original Proposition10 and complete author proof plus inner-product Corollaries11--12, and full Section5 Proposition18/Lemma19/Theorem9 proofs retained. The projection lemma includes all averaging, Markov, coefficient, orthonormal-basis and constant13 calculations and original footnote. Main proof includes complete permutation extension, every completion/agreement argument, full two-slice inner-product expansion, and low-degree projection/spectral concentration conclusion. All36 original bibliography entries retained. Main arbitrary separate preprocessing, short Bob output n+k, constant depth/size bounds and IP implications are unchanged. Original printed wording/notation is preserved rather than mathematically repaired: Theorem9 right one-to-one premise versus proofSTEP1 left wording; STEP2 printed vector codomain; printed y\_(n+i) positions; Proposition10 displayed expectation inequality; original k=0 range/final ratios. No added hypotheses or independent correctness claim. The author-standard external LMNT/entropy statements retained without recursive reproof. The2024 journal57-page version was retrieved but its added correlation results are not mixed into this item; exact published CCC proof is complete. Conditional learning/cryptographic/rounded-IP applications are outside this item and not separately counted. No source-native TeX was found in bounded lookup, so the licensed formal PDF/text was faithfully transcribed. No figures required. Root handles compilation/database/publication.\par\medskip
\noindent\textit{Allowed standard prerequisites (editorial).}\begin{itemize}
\item Tal LMN spectral concentration improvement. Original Theorem7[24,30] exact statement retained; Proposition18 full author deduction included.
\item Binary-entropy bound. Original Lemma8[21,Lemma3.6] exact statement retained.
\item Fourier characters and orthogonal projections. Original Section2 inner product/orthonormal basis and Section5 projection notation retained; standard linear algebra not recursively reproved.
\end{itemize}\medskip
\setcounter{nttheorem}{8}
\begin{nttheorem}[Author Theorem9: Main Theorem]
Let $f:\{-1,1\}^n\times\{-1,1\}^n\to\{-1,1\}$ be a Boolean function, let $0\leq k\leq n/2-1$, and let $0\leq t\leq n+k$ be an integer. Denote $f_x(y)\triangleq f(x,y)$. Suppose the following hold:
\begin{itemize}
\item $f$ satisfies the right one-to-one condition.
\item There exists a subset $X\subseteq\{-1,1\}^n$ of size $|X|\geq13\cdot2^{2(k+1)}\cdot\binom{n+k}{\leq t}$ such that
\[\mathbb E_{x\ne x'\sim X}[\langle f_x,f_{x'}\rangle^2]\leq\frac{2^{2k}}{36|X|^2}.\]
\item $f$ admits an SM protocol $P=(A,B,C)$ such that $B:\{-1,1\}^n\to\{-1,1\}^{n+k}$ and $C$ is an unbounded fan-in circuit of depth $h$ and size $M$.
\end{itemize}
Then:
\[M\geq2^{\Omega_h([t/k]^{1/(h-1)})}.\]
\end{nttheorem}

\section*{2 Preliminaries}
\subsection*{2.1 Definitions and notation}
\paragraph{Simultaneous messages protocols}
A (two-party) simultaneous messages (SM) protocol consists of two players, which we refer to as Alice and Bob, and a referee, which we refer to as Carol, that together compute a function. Alice and Bob each send a message, which is based on the input, to Carol, who then computes a function of the two messages received. Formally, we have the following definitions:
\setcounter{nttheorem}{2}
\begin{ntdefinition}[Author Definition3: Simultaneous messages protocols]
Let $X,Y,\widehat X,\widehat Y,Z$ be finite nonempty sets. A simultaneous messages protocol (shortly, SM protocol or SMP) $P$ is a triplet of functions $(A,B,C)$, where $A:X\to\widehat X$, $B:Y\to\widehat Y$, and $C:\widehat X\times\widehat Y\to Z$. We call $C$ the referee function.
\end{ntdefinition}
\begin{ntdefinition}[Author Definition4: SM protocol admittance]
Let $f:X\times Y\to Z$ be a function. We say that $f$ admits an SM protocol $(A,B,C)$ if $f(x,y)=C(A(x),B(y))$ for every $(x,y)\in X\times Y$. In that case, we also say that $(A,B,C)$ computes $f$.
\end{ntdefinition}
\paragraph{Inner product space of Boolean functions}
For the purpose of utilizing Fourier analysis, we will consider the inner product space of all functions $\{-1,1\}^n\to\mathbb R$ with the following inner product:
\[\langle f,g\rangle=\mathbb E_{x\sim\{-1,1\}^n}[f(x)g(x)]=\frac1{2^n}\sum_{x\in\{-1,1\}^n}f(x)\cdot g(x).\]
It is a known fact that the aforementioned inner product space has as orthonormal basis the set of all parity functions $\{\chi_S\}_{S\subseteq[n]}$, defined by $\chi_S(x)=\prod_{i\in S}x_i$.
\paragraph{The Inner Product function}
The inner product modulo2 function $\operatorname{IP}:\{0,1\}^n\times\{0,1\}^n\to\{0,1\}$ is defined by
\[\operatorname{IP}(x,y)=\sum_{i=1}^nx_iy_i\pmod2.\]
Note that each $x$, when fixed, corresponds to a parity function on a subset $S$ of $y$'s coordinates, a subset which is determined by $x$. This correspondence is actually a bijection mapping elements of $\{0,1\}^n$ to subsets of $2^{[n]}$. Thus, when we switch to $\pm1$ notation, we can identify each $x\in\{-1,1\}^n$ with the parity function $\chi_{S(x)}(y)$ that results when fixing $x$ in $\operatorname{IP}(x,y)$.
\paragraph{One-to-one condition}
Some functions have a certain property, formally defined below, rendering them harder to compute with preprocessing. Restricting attention to these functions seems to help in obtaining lower bounds.
\begin{ntdefinition}[Author Definition5: One-to-one condition]
Let $f:X\times Y\to Z$ be a function. We say that $f$ satisfies the left one-to-one condition if for every $x\ne x'\in X$ there exists $y\in Y$ such that $f(x,y)\ne f(x',y)$. Similarly, we say that $f$ satisfies the right one-to-one condition if for every $y\ne y'\in Y$ there exists $x\in X$ such that $f(x,y)\ne f(x,y')$. Finally, we say that $f$ satisfies the one-to-one condition if $f$ satisfies both the left and right one-to-one conditions.
\end{ntdefinition}
\begin{ntproposition}[Author Proposition6: One-to-one preprocessing]
Let $f:X\times Y\to Z$ be a function. Then:
\begin{itemize}
\item If $f$ satisfies the left one-to-one condition, then for every SM protocol $(A,B,C)$ that $f$ admits, $A$ computes a one-to-one mapping.
\item If $f$ satisfies the right one-to-one condition, then for every SM protocol $(A,B,C)$ that $f$ admits, $B$ computes a one-to-one mapping.
\end{itemize}
\end{ntproposition}
\subsection*{2.2 Known facts}
The following are known facts we will need later.
\begin{nttheorem}[Author Theorem7: Tal's LMN improvement (LMNT), {\cite{ref24,ref30}}]
Let $f$ be a Boolean function with $n$ variables computable by an unbounded fan-in circuit of depth $h$ and size $M$, and let $t$ be any integer. Then,
\[\|f^{\geq t}\|^2\leq2\cdot2^{-t/O_h(\log M)^{h-1}}.\]
\end{nttheorem}
\begin{ntlemma}[Author Lemma8, Lemma3.6 in {\cite{ref21}}]
Let $H:[0,1]\to\mathbb R$ be the binary entropy function defined by $H(p)=-p\log p-(1-p)\log(1-p)$. Then, for any $0<a\leq1/2$ and $n\in\mathbb N$,
\[\binom n{\leq an}\leq2^{H(a)n}.\]
\end{ntlemma}

\section*{3 Main results: Inner-product consequences}
Our main technical result, proved in Section5, states that a ``large'' collection of functions that are ``close'' to being orthonormal, is computationally hard for SM protocols in which the referee is an unbounded fan-in circuit of constant depth, and one player is limited to ``short'' preprocessing output length.
One may wonder about the necessity of satisfying the one-to-one condition. While it may still be the case that the same result (or even better) follows without this assumption, if we could obtain it that way, then we could easily exhibit a function with stronger lower bounds. As an example, consider taking the inner product function and computing it only on a prefix of the input while ignoring other bits--relying on the main theorem's consequence, we could extend it to an arbitrary $\Omega(n)$ lower bound on the preprocessing output length.
The (somewhat cumbersome) second requirement on a function $f:\{-1,1\}^n\times\{-1,1\}^n\to\{-1,1\}$ specified in Theorem9 can be replaced by a slightly stronger requirement, yet one involving the more familiar measure borrowed from the study of learning by statistical queries~\cite{ref23}. Define the statistical query dimension of $f$ with respect to the uniform distribution to be the size of the largest set $D\subseteq\{-1,1\}^n$ such that $|\langle f_x,f_{x'}\rangle|\leq1/|D|$ for every $x\ne x'\in D$. More details on statistical query dimension can be found in~\cite{ref7}.
A simple consequence of Theorem9 is Theorem1. Here we state a more general version that refers to statistical query dimension. At a high level, the theorem says that computing with preprocessing a function having exponential statistical query dimension remains as hard for $\operatorname{AC}^0$ as without, given that one player is limited to output a string whose length stretches the input length by an additive sublinear term:
\setcounter{nttheorem}{9}
\begin{ntproposition}[Author Proposition10: Formal version of Theorem1]
Let $k=o(n)$, and suppose that a function $f:\{-1,1\}^n\times\{-1,1\}^n\to\{-1,1\}$ satisfies the one-to-one condition and has statistical query dimension of $2^{\Omega(n)}$. If $f$ admits an SM protocol $(A,B,C)$ such that $B:\{0,1\}^n\to\{0,1\}^{n+k}$ and $C$ is an unbounded fan-in circuit of depth $h$ and size $M$, then:
\[M\geq2^{\Omega_h([n/k]^{1/(h-1)})}.\]
\end{ntproposition}
\begin{proof}
Let $D\subseteq\{-1,1\}^n$ be a set of size $2^{\Omega(n)}$ such that $|\langle f_x,f_{x'}\rangle|\leq1/|D|$ for every $x\ne x'\in D$. One can easily find an $0<\alpha\leq1/2$ for which $H(\alpha)$ is small enough, such that setting $t=\alpha(n+k)$ gives
\[13\cdot2^{2(k+1)}\cdot\binom{n+k}{\leq\alpha(n+k)}\ \underset{\text{Lemma8}}\leq\ 13\cdot2^{2(k+1)}\cdot2^{H(\alpha)(n+k)}\leq|D|/6.\]
Now, let $X\subseteq D$ of size $|X|=13\cdot2^{2(k+1)}\cdot\binom{n+k}{\leq t}$. We have:
\begin{align*}\mathbb E_{x\ne x'\sim X}[\langle f_x,f_{x'}\rangle^2]&\leq\mathbb E_{x\ne x'\sim D}[\langle f_x,f_{x'}\rangle^2]\leq\left(\mathbb E_{x\ne x'\sim D}[\langle f_x,f_{x'}\rangle]\right)^2\\&\leq\frac1{|D|^2}\leq\frac1{36|X|^2}\leq\frac{2^{2k}}{36|X|^2}.\end{align*}
Thus, by Theorem9,
\[M\geq2^{\Omega_h([t/k]^{1/(h-1)})}=2^{\Omega_h([\alpha(n+k)/k]^{1/(h-1)})}=2^{\Omega_h([n/k]^{1/(h-1)})}.\]
\end{proof}
Since $\operatorname{IP}$ satisfies the one-to-one condition and has the largest possible statistical query dimension of $2^n$, we get the following corollary.
\begin{ntcorollary}[Author Corollary11]
Let $k\leq n^\alpha$ for some $0\leq\alpha<1$, and suppose that $\operatorname{IP}$ admits an SM protocol $(A,B,C)$ such that $B:\{0,1\}^n\to\{0,1\}^{n+k}$ and $C$ is an unbounded fan-in circuit of depth $h$ and size $M$. Then:
\[M\geq2^{\Omega_h(n^{(1-\alpha)/(h-1)})}.\]
\end{ntcorollary}
\begin{ntcorollary}[Author Corollary12]
Let $k\leq n/\log^\beta n$ for every $\beta>0$ (for large enough $n$), and suppose that $\operatorname{IP}$ admits an SM protocol $(A,B,C)$ such that $B:\{0,1\}^n\to\{0,1\}^{n+k}$ and $C$ is an unbounded fan-in circuit of depth $h$ and size $M$. Then:
\[M\geq2^{\Omega_h(\log^c n)}\qquad\text{for every }c>0.\]
\end{ntcorollary}

\section*{5 Proof of Main Theorem}
The following two results will be needed for proving the main theorem, Theorem9.
\setcounter{nttheorem}{17}
\begin{ntproposition}[Author Proposition18: High-degree spectral concentration bound]
Let $f:\{-1,1\}^n\to\{-1,1\}$ be a Boolean function, and let $0\leq\epsilon\leq1/2$. Then, for every integer $0\leq t\leq n$ such that $\|f^{\leq t}\|\leq\epsilon$, if an unbounded fan-in circuit of depth $h$ and size $M$ agrees with $f$ on at least $1/2+\epsilon$ fraction of inputs, then
\[M\geq2^{\Omega_h([t/(1-2\log\epsilon)]^{1/(h-1)})}.\]
\end{ntproposition}
\begin{proof}
Let $0\leq t\leq n$ be an integer such that $\|f^{\leq t}\|\leq\epsilon$, and suppose that an unbounded fan-in circuit of depth $h$ and size $M$ computes a function $F$ that agrees with $f$ on at least $1/2+\epsilon$ fraction of inputs.
On one hand, we have
\[\langle F,f\rangle=2\Pr[F=f]-1\ \underset{\text{assumption}}\geq\ 2(1/2+\epsilon)-1=2\epsilon.\]
On the other hand, we have
\begin{align*}\langle F,f\rangle&=\langle F^{\leq t},f^{\leq t}\rangle+\langle F^{>t},f^{>t}\rangle\\&\underset{\text{Cauchy--Schwarz}}\leq\|f^{\leq t}\|+\|F^{>t}\|\ \underset{\text{LMNT}}\leq\ \epsilon+\sqrt{2\cdot2^{-t/O_h(\log M)^{h-1}}}.\end{align*}
Thus,
\[2\epsilon\leq\epsilon+\sqrt{2\cdot2^{-t/O_h(\log M)^{h-1}}}\quad\Longrightarrow\quad M\geq2^{\Omega_h([t/(1-2\log\epsilon)]^{1/(h-1)})}.\]
\end{proof}
In what follows, we will use the following notation:
\begin{itemize}
\item For a set $X$, we write $i\ne j\sim X$ to mean that $(i,j)$ is chosen uniformly at random from the set $\{(i,j)\in X\times X:i\ne j\}$.
\item Given an inner product space $V$, a subspace $U\leq V$, and a vector $v\in V$, we denote the projection of $v$ onto $U$ by $\operatorname{proj}_U(v)$.
\end{itemize}

\setcounter{nttheorem}{18}
\begin{ntlemma}[Author Lemma19: The Projection Lemma]
Let $V$ be an inner product space over $\mathbb R$. Let $\{v_i\}_{i\in X}\subseteq V$ be a set of unit vectors indexed by $X$, and suppose that
\[\mathbb E_{i\ne j\sim X}[\langle v_i,v_j\rangle^2]\leq\frac1{36|X|^2}.\]
Then, for every subspace $U\leq V$, there exists $i\in X$ such that $\|\operatorname{proj}_U(v_i)\|^2=O(\dim U/|X|)$.
\end{ntlemma}
\begin{proof}
Let $U\leq V$ be a subspace, and denote $D\triangleq\dim U$.
By Cauchy--Schwartz,
\[\mathbb E_{i\ne j\sim X}[|\langle v_i,v_j\rangle|]\leq\left(\mathbb E_{i\ne j\sim X}[\langle v_i,v_j\rangle^2]\right)^{1/2}\leq\frac1{6|X|}.\]
By Markov's inequality,
\[\Pr_{i\sim X}\left[\mathbb E_{j\sim X\setminus\{i\}}[\langle v_i,v_j\rangle^2]>\frac1{12|X|^2}\right]\leq12|X|^2\cdot\mathbb E_{i\ne j}[\langle v_i,v_j\rangle^2]\leq\frac13,\]
and similarly,
\[\Pr_{i\sim X}\left[\mathbb E_{j\sim X\setminus\{i\}}[|\langle v_i,v_j\rangle|]>\frac1{2|X|}\right]\leq2|X|\cdot\mathbb E_{i\ne j}[|\langle v_i,v_j\rangle|]\leq\frac13,\]
which implies that at least $1/3$ of the indices $i\in X$ satisfy
\[\mathbb E_{j\sim X\setminus\{i\}}[\langle v_i,v_j\rangle^2]\leq\frac1{12|X|^2}\quad\text{and}\quad\mathbb E_{j\sim X\setminus\{i\}}[|\langle v_i,v_j\rangle|]\leq\frac1{2|X|},\]
or equivalently,
\begin{equation}\sum_{j\in X\setminus\{i\}}\langle v_i,v_j\rangle^2\leq\frac1{12|X|}\quad\text{and}\quad\sum_{j\in X\setminus\{i\}}|\langle v_i,v_j\rangle|\leq\frac12.\tag{1}\end{equation}
Put these indices in a set $Y$, and let $W\triangleq\operatorname{span}(\{v_i:i\in Y\})$.
Let $w\in W$ such that $\|w\|\leq1$, and write $w=\sum_{i\in Y}c_iv_i$ with $c_i\in\mathbb R$. Then for $i\in Y$,
\[\langle w,v_i\rangle=\sum_{j\in Y}c_j\langle v_i,v_j\rangle=c_i+\sum_{j\in Y\setminus\{i\}}c_j\langle v_i,v_j\rangle.\]
Multiply this by $c_i$, and sum over all $i\in Y$ to obtain
\[1\geq\|w\|^2=\sum_{i\in Y}c_i^2+\sum_{i\ne j}c_ic_j\langle v_i,v_j\rangle.\]
Since $2|c_ic_j|\leq c_i^2+c_j^2$, it follows that
\begin{align*}1&\geq\sum_{i\in Y}c_i^2-\frac12\sum_{i\ne j}(c_i^2+c_j^2)|\langle v_i,v_j\rangle|\\&=\sum_{i\in Y}c_i^2\left(1-\sum_{j\in Y\setminus\{i\}}|\langle v_i,v_j\rangle|\right)\ \underset{\text{Eq.(1)}}\geq\ \frac12\sum_{i\in Y}c_i^2,\end{align*}
which implies $\sum_{i\in Y}c_i^2\leq2$.\footnote{Note that the argument implies that the vectors $\{v_i\}_{i\in Y}$ are linearly independent; otherwise, we can find representations of $w$ for which $\sum c_i^2$ is arbitrary large.} Since $(a+b)^2\leq2a^2+2b^2$, for every $i\in Y$ we have
\begin{align*}\langle w,v_i\rangle^2&\leq2c_i^2+2\left(\sum_{j\in Y\setminus\{i\}}c_j\langle v_i,v_j\rangle\right)^2\\&\underset{\text{Cauchy--Schwarz}}\leq2c_i^2+2\sum_{j\in Y\setminus\{i\}}c_j^2\cdot\sum_{j\in Y\setminus\{i\}}\langle v_i,v_j\rangle^2\\&\leq2c_i^2+4\sum_{j\in Y\setminus\{i\}}\langle v_i,v_j\rangle^2\ \underset{\text{Eq.(1)}}\leq\ 2c_i^2+\frac1{3|X|}.\end{align*}
Taking expectation over $i\in Y$, we deduce
\begin{align*}\mathbb E_{i\sim Y}[\langle w,v_i\rangle^2]&\leq\mathbb E_{i\sim Y}\left[2c_i^2+\frac1{3|X|}\right]=\frac2{|Y|}\sum_{i\in Y}c_i^2+\frac1{3|X|}\\&\leq\frac4{|Y|}+\frac1{3|X|}\ \underset{|Y|\geq|X|/3}\leq\ \frac{12}{|X|}+\frac1{3|X|}\leq\frac{13}{|X|}.\end{align*}
Now let $u_1,\ldots,u_D$ be an orthonormal basis for $U$, and for every $k\in[D]$, let $w_k\in W$ be the projection of $u_k$ onto $W$ (notice that $\|w_k\|\leq1$). We have
\begin{align*}\mathbb E_{i\sim Y}[\|\operatorname{proj}_U(v_i)\|^2]&=\mathbb E_{i\sim Y}\left[\sum_{k\in[D]}\langle v_i,u_k\rangle^2\right]=\mathbb E_{i\sim Y}\left[\sum_{k\in[D]}\langle v_i,w_k\rangle^2\right]\\&=\sum_{k\in[D]}\mathbb E_{i\sim Y}[\langle v_i,w_k\rangle^2]\leq\frac{13D}{|X|},\end{align*}
which implies there exists $i\in Y$ such that $\|\operatorname{proj}_U(v_i)\|^2\leq13D/|X|=O(D/|X|)$, as desired.
\end{proof}
We can now prove our main theorem.

\begin{proof}[Author proof of Theorem9]
The proof follows several steps.
\paragraph{STEP1}
Since $f$ satisfies the left one-to-one condition, by Proposition6, $B$ computes a one-to-one mapping; hence, we can extend it to a permutation $\tau:\{-1,1\}^{n+k}\to\{-1,1\}^{n+k}$ as follows:
\[\tau(y_1,\ldots,y_{n+k})=\begin{cases}B(y_1,\ldots,y_n)&\text{if }y_{n+1}=\cdots=y_{n+k}=1,\\\text{arbitrary choice}&\text{otherwise},\end{cases}\]
where by arbitrary choice we mean one of the $(2^{n+k}-2^n)!$ possible ways of completing the definition so as to yield a permutation. Define $\sigma=\tau^{-1}$ and note that $\sigma$ is a permutation as well.
\paragraph{STEP2}
For every $x\in\{-1,1\}^n$ and $R\subseteq\{n+1,\ldots,n+k\}$, define $f_x^R:\{-1,1\}^{n+k}\to\{-1,1\}^{n+k}$ by
\[f_x^R(y_1,\ldots,y_{n+k})=\begin{cases}f_x(y_1,\ldots,y_n)&\text{if }y_{n+1}=\cdots=y_{n+k}=1,\\\chi_{S(x)}(y_1,\ldots,y_n)\cdot\chi_R(y_{n+1},\ldots,y_{n+k})&\text{otherwise}.\end{cases}\]
What can we say about these functions?
\begin{itemize}
\item Fix $x\in\{-1,1\}^n$, and denote by $C_x$ the circuit obtained from $C$ when Alice is given $x$ as input. Now, consider $y\in\{-1,1\}^{n+k}$.
\begin{itemize}
\item If $y_{n+1}=\cdots=y_{n+k}=1$, then $f_x^R(y)=f_x(y)$ by definition; hence, by the correctness of $P$ and the definition of $\sigma$, we have that $f_x^R$ agrees with $C_x\circ\sigma^{-1}$ on all such $y$'s.
\item Otherwise, let $i\in\{n+1,\ldots,n+k\}$ such that $y_i=-1$. For every $R\subseteq\{n+1,\ldots,n+k\}$ that contains $i$, we have that $C_x\circ\sigma^{-1}$ agrees with exactly one of $f_x^R,f_x^{R\setminus\{i\}}$ on the input $(y_1,\ldots,y_{n+i-1},-1,y_{n+i+1},\ldots,y_{n+k})$; thus, for exactly half the subsets $R\subseteq\{n+1,\ldots,n+k\}$, $f_x^R$ agrees with $C_x\circ\sigma^{-1}$ on $y$. Therefore,
\[\Pr_{R\sim2^{[n+k]\setminus[n]}}[f_x^R(y)=C_x(\sigma^{-1}(y))]=\frac12.\]
This holds for any $y$ such that $(y_{n+1},\ldots,y_{n+k})\ne(1,\ldots,1)$; hence,
\[\underset{\substack{y\sim\{-1,1\}^{n+k}\\\exists j\in[k]:y_{n+j}=-1}}{\mathbb E}\left[\Pr_{R\sim2^{[n+k]\setminus[n]}}[f_x^R(y)=C_x(\sigma^{-1}(y))]\right]=\frac12,\]
which implies there exists $R(x)\subseteq[n+k]\setminus[n]$ such that
\[\Pr_{\substack{y\sim\{-1,1\}^{n+k}\\\exists j\in[k]:y_{n+j}=-1}}[f_x^{R(x)}(y)=C_x(\sigma^{-1}(y))]\geq\frac12.\]
\end{itemize}
It follows that the fraction of inputs on which $f_x^{R(x)}$ and $C_x\circ\sigma^{-1}$ agree is at least
\[\frac{2^n+(1/2)\cdot(2^{n+k}-2^n)}{2^{n+k}}=\frac12+\frac1{2^{k+1}},\]
which is also the fraction of inputs on which $F_x^{R(x)}\triangleq f_x^{R(x)}\circ\sigma$ and $C_x$ agree.
\item The second thing we observe is that for any $x\ne x'\in X$,
\begin{align*}\langle F_x^{R(x)},F_{x'}^{R(x')}\rangle&=\langle f_x^{R(x)}\circ\sigma,f_{x'}^{R(x')}\circ\sigma\rangle=\langle f_x^{R(x)},f_{x'}^{R(x')}\rangle\\&=\mathbb E_{y\sim\{-1,1\}^{n+k}}[f_x^{R(x)}(y)\cdot f_{x'}^{R(x')}(y)]\\&=\mathbb E_{y\sim\{-1,1\}^n}[f_x(y)\cdot f_{x'}(y)]\cdot2^{-k}\\&\quad+\sum_{\substack{z\in\{-1,1\}^k\\\exists j\in[k]:z_j=-1}}\chi_{R(x)}(z)\cdot\chi_{R(x')}(z)\cdot\mathbb E_{y\sim\{-1,1\}^n}[\chi_{S(x)}(y)\cdot\chi_{S(x')}(y)]\cdot2^{-k}\\&=\langle f_x,f_{x'}\rangle\cdot2^{-k}\\&\quad+\sum_{\substack{z\in\{-1,1\}^k\\\exists j\in[k]:z_j=-1}}\chi_{R(x)}(z)\cdot\chi_{R(x')}(z)\cdot\underbrace{\langle\chi_{S(x)},\chi_{S(x')}\rangle}_{0}\cdot2^{-k}\\&=\langle f_x,f_{x'}\rangle\cdot2^{-k},\end{align*}
which implies
\begin{align*}\mathbb E_{x\ne x'\sim X}[\langle F_x^{R(x)},F_{x'}^{R(x')}\rangle^2]&\leq2^{-2k}\cdot\mathbb E_{x\ne x'\sim X}[\langle f_x,f_{x'}\rangle^2]\\&\underset{\text{assumption}}\leq2^{-2k}\cdot\frac{2^{2k}}{36|X|^2}=\frac1{36|X|^2}.\end{align*}
\end{itemize}
\paragraph{STEP3}
Let $V$ be the inner product space of all functions $\{-1,1\}^{n+k}\to\mathbb R$, and let $U\leq V$ be the subspace of all functions of degree up to $t$, which is spanned by $\{\chi_T\}_{|T|\leq t}$ and has dimension $\dim U=\binom{n+k}{\leq t}$. By the Projection Lemma, there exists $x^*\in X$ such that
\[\|(F_{x^*}^{R(x^*)})^{\leq t}\|^2\leq\frac{13\binom{n+k}{\leq t}}{|X|}\ \underset{\text{assumption}}\leq\ \frac1{2^{2(k+1)}}\quad\Longrightarrow\quad\|(F_{x^*}^{R(x^*)})^{\leq t}\|\leq\frac1{2^{k+1}}.\]
Since $\|(F_{x^*}^{R(x^*)})^{\leq t}\|\leq1/2^{k+1}$ and $C_{x^*}$ is an unbounded fan-in circuit of depth $h$ and size $\leq M$ that agrees with $F_{x^*}^{R(x^*)}$ on at least $1/2+1/2^{k+1}$ fraction of inputs, by Proposition18,
\[M\geq2^{\Omega_h([t/(1-2\log2^{-(k+1)})]^{1/(h-1)})}=2^{\Omega_h([t/(2k+3)]^{1/(h-1)})}=2^{\Omega_h([t/k]^{1/(h-1)})}.\]
\end{proof}

\begin{thebibliography}{36}
\bibitem{ref1}Miklós Ajtai. Σ11 -formulae on finite structures. Ann. Pure Appl. Logic, 24(1):1–48, 1983. doi:10.1016/0168-0072(83)90038-6.
\bibitem{ref2}Adi Akavia, Andrej Bogdanov, Siyao Guo, Akshay Kamath, and Alon Rosen. Candidate weak pseudorandom functions in AC0 ◦ MOD2 . In Innovations in Theoretical Computer Science, ITCS’14, Princeton, NJ, USA, January 12-14, 2014, pages 251–260, 2014. doi: 10.1145/2554797.2554821.
\bibitem{ref3}Josh Alman and R. Ryan Williams. Probabilistic rank and matrix rigidity. In Proceedings of the 49th Annual ACM SIGACT Symposium on Theory of Computing, STOC 2017, Montreal, QC, Canada, June 19-23, 2017, pages 641–652, 2017. doi:10.1145/3055399.3055484.
\bibitem{ref4}László Babai, Peter Frankl, and Janos Simon. Complexity classes in communication complexity theory (preliminary version). In 27th Annual Symposium on Foundations of Computer Science, Toronto, Canada, 27-29 October 1986, pages 337–347, 1986. doi:10.1109/SFCS.1986.15.
\bibitem{ref5}László Babai, Anna Gál, Peter G. Kimmel, and Satyanarayana V. Lokam. Communication complexity of simultaneous messages. SIAM J. Comput., 33(1):137–166, 2003. doi:10.1137/ S0097539700375944.
\bibitem{ref6}Abhishek Banerjee, Chris Peikert, and Alon Rosen. Pseudorandom functions and lattices. In Advances in Cryptology - EUROCRYPT 2012 - 31st Annual International Conference on the Theory and Applications of Cryptographic Techniques, Cambridge, UK, April 15-19, 2012. Proceedings, pages 719–737, 2012. doi:10.1007/978-3-642-29011-4\_42.
\bibitem{ref7}Avrim Blum, Merrick L. Furst, Jeffrey C. Jackson, Michael J. Kearns, Yishay Mansour, and Steven Rudich. Weakly learning DNF and characterizing statistical query learning using fourier analysis. In Proceedings of the Twenty-Sixth Annual ACM Symposium on Theory of Computing, 23-25 May 1994, Montréal, Québec, Canada, pages 253–262. ACM, 1994. doi:10.1145/195058.195147.
\bibitem{ref8}Andrej Bogdanov, Yuval Ishai, and Akshayaram Srinivasan. Unconditionally secure compu- tation against low-complexity leakage. In Advances in Cryptology - CRYPTO 2019 - 39th Annual International Cryptology Conference, Santa Barbara, CA, USA, August 18-22, 2019, Proceedings, Part II, volume 11693 of Lecture Notes in Computer Science, pages 387–416, 2019. doi:10.1007/978-3-030-26951-7\_14.
\bibitem{ref9}Andrej Bogdanov and Alon Rosen. Pseudorandom functions: Three decades later. In Tutorials on the Foundations of Cryptography, pages 79–158. Springer, Cham, 2017. doi: 10.1007/978-3-319-57048-8\_3.
\bibitem{ref10}Dan Boneh, Yuval Ishai, Alain Passelègue, Amit Sahai, and David J. Wu. Exploring crypto dark matter: - new simple PRF candidates and their applications. In Theory of Cryptography - 16th International Conference, TCC 2018, Panaji, India, November 11-14, 2018, Proceedings, Part II, volume 11240 of Lecture Notes in Computer Science, pages 699–729, 2018. doi: 10.1007/978-3-030-03810-6\_25.
\bibitem{ref11}Arkadev Chattopadhyay and Rahul Santhanam. Lower bounds on interactive compressibility by constant-depth circuits. In 53rd Annual IEEE Symposium on Foundations of Computer Science, FOCS 2012, New Brunswick, NJ, USA, October 20-23, 2012, pages 619–628, 2012. doi:10.1109/FOCS.2012.74.
\bibitem{ref12}Mahdi Cheraghchi, Elena Grigorescu, Brendan Juba, Karl Wimmer, and Ning Xie. AC0 ◦MOD2 lower bounds for the boolean inner product. In 43rd International Colloquium on Automata, Languages, and Programming, ICALP 2016, July 11-15, 2016, Rome, Italy, volume 55 of LIPIcs, pages 35:1–35:14, 2016. doi:10.4230/LIPIcs.ICALP.2016.35.
\bibitem{ref13}Amit Daniely, Nati Linial, and Shai Shalev-Shwartz. From average case complexity to improper learning complexity. In Symposium on Theory of Computing, STOC 2014, New York, NY, USA, May 31 - June 03, 2014, pages 441–448, 2014. doi:10.1145/2591796.2591820.
\bibitem{ref14}Ning Ding, Yanli Ren, and Dawu Gu. PAC learning depth-3 AC0 circuits of bounded top fanin. In International Conference on Algorithmic Learning Theory, ALT 2017, 15-17 October 2017, Kyoto University, Kyoto, Japan, pages 667–680, 2017. URL: http://proceedings.mlr. press/v76/ding17a.html.
\bibitem{ref15}Bella Dubrov and Yuval Ishai. On the randomness complexity of efficient sampling. In Proceedings of the 38th Annual ACM Symposium on Theory of Computing, Seattle, WA, USA, May 21-23, 2006, pages 711–720, 2006. doi:10.1145/1132516.1132615.
\bibitem{ref16}Zeev Dvir and Benjamin Edelman. Matrix rigidity and the croot-lev-pach lemma. CoRR, abs/1708.01646, 2017. arXiv:1708.01646.
\bibitem{ref17}Merrick L. Furst, James B. Saxe, and Michael Sipser. Parity, circuits, and the polynomial- time hierarchy. In 22nd Annual Symposium on Foundations of Computer Science, Nashville, Tennessee, USA, 28-30 October 1981, pages 260–270, 1981. doi:10.1109/SFCS.1981.35.
\bibitem{ref18}Edgar N. Gilbert. A comparison of signalling alphabets. The Bell system technical journal, 31(3):504–522, 1952.
\bibitem{ref19}Oded Goldreich, Shafi Goldwasser, and Silvio Micali. How to construct random functions. J. ACM, 33(4):792–807, 1986. doi:10.1145/6490.6503.
\bibitem{ref20}Mika Göös, Toniann Pitassi, and Thomas Watson. The landscape of communication complexity classes. In 43rd International Colloquium on Automata, Languages, and Programming, ICALP 2016, July 11-15, 2016, Rome, Italy, volume 55 of LIPIcs, pages 86:1–86:15, 2016. doi: 10.4230/LIPIcs.ICALP.2016.86.
\bibitem{ref21}Robert M. Gray. Entropy and information theory. Springer Science \& Business Media, 2011.
\bibitem{ref22}Johan Håstad. Almost optimal lower bounds for small depth circuits. In Proceedings of the 18th Annual ACM Symposium on Theory of Computing, May 28-30, 1986, Berkeley, California, USA, pages 6–20, 1986. doi:10.1145/12130.12132.
\bibitem{ref23}Michael J. Kearns. Efficient noise-tolerant learning from statistical queries. In Proceedings of the Twenty-Fifth Annual ACM Symposium on Theory of Computing, May 16-18, 1993, San Diego, CA, USA, pages 392–401. ACM, 1993. doi:10.1145/167088.167200.
\bibitem{ref24}Nathan Linial, Yishay Mansour, and Noam Nisan. Constant depth circuits, fourier transform, and learnability. In 30th Annual Symposium on Foundations of Computer Science, Research Triangle Park, North Carolina, USA, 30 October - 1 November 1989, pages 574–579, 1989. doi:10.1109/SFCS.1989.63537.
\bibitem{ref25}Pavel Pudlák, Vojtech Rödl, and Petr Savický. Graph complexity. Acta Inf., 25(5):515–535, 1988. doi:10.1007/BF00279952.
\bibitem{ref26}Alexander A. Razborov. On rigid matrices. preprint, 1989.
\bibitem{ref27}Guy N. Rothblum. How to compute under AC0 leakage without secure hardware. In Advances in Cryptology - CRYPTO 2012 - 32nd Annual Cryptology Conference, Santa Barbara, CA, USA, August 19-23, 2012. Proceedings, volume 7417 of Lecture Notes in Computer Science, pages 552–569, 2012. doi:10.1007/978-3-642-32009-5\_32.
\bibitem{ref28}Rocco A. Servedio and Emanuele Viola. On a special case of rigidity. Electronic Colloquium on Computational Complexity (ECCC), 19:144, 2012. URL: http://eccc.hpi-web.de/report/ 2012/144.
\bibitem{ref29}Avishay Tal. The bipartite formula complexity of inner-product is quadratic. Electronic Colloquium on Computational Complexity (ECCC), 23:181, 2016. URL: http://eccc.hpi-web. de/report/2016/181.
\bibitem{ref30}Avishay Tal. Tight bounds on the fourier spectrum of AC0 . In 32nd Computational Complexity Conference, CCC 2017, July 6-9, 2017, Riga, Latvia, volume 79 of LIPIcs, pages 15:1–15:31, 2017. doi:10.4230/LIPIcs.CCC.2017.15.
\bibitem{ref31}Salil P. Vadhan. On learning vs. refutation. In Proceedings of the 30th Conference on Learning Theory, COLT 2017, Amsterdam, The Netherlands, 7-10 July 2017, pages 1835–1848, 2017. URL: http://proceedings.mlr.press/v65/vadhan17a.html.
\bibitem{ref32}Leslie G. Valiant. A theory of the learnable. In Proceedings of the 16th Annual ACM Symposium on Theory of Computing, April 30 - May 2, 1984, Washington, DC, USA, pages 436–445, 1984. doi:10.1145/800057.808710.
\bibitem{ref33}R. R. Varshamov. Estimate of the number of signals in error correcting codes. Docklady Akad. Nauk, SSSR, 117:739–741, 1957.
\bibitem{ref34}Emanuele Viola. Extractors for circuit sources. In IEEE 52nd Annual Symposium on Founda- tions of Computer Science, FOCS 2011, Palm Springs, CA, USA, October 22-25, 2011, pages 220–229, 2011. doi:10.1109/FOCS.2011.20.
\bibitem{ref35}Henning Wunderlich. On a theorem of razborov. Computational Complexity, 21(3):431–477, 2012. doi:10.1007/s00037-011-0021-5.
\bibitem{ref36}Andrew Chi-Chih Yao. Some complexity questions related to distributive computing (prelimin- ary report). In Proceedings of the 11h Annual ACM Symposium on Theory of Computing, April 30 - May 2, 1979, Atlanta, Georgia, USA, pages 209–213, 1979. doi:10.1145/800135.804414.
\end{thebibliography}
\end{document}
